Continuation of Reviewer VCqo answer

---

Rebuttal (Reviewer 2 --- full answers, matching \texttt{answers.tex})

We thank the reviewer for the careful and substantive feedback. What follows is the complete response text we stand behind (we do not compress here).

Q1 (Novelty, freezing, applicability, Theorem~4.1)

We agree with the spirit of this comment: the proof template is recognizably related to Nguyen \& Pham (2023), and the two-point theorem is a minimal rigorous instance rather than a full theory of representation learning on natural data. Where we respectfully disagree is the framing that freezing is only a limitation. For our goals, freezing the mixing matrices is simultaneously a modeling choice and a technical enabler. It is what makes the deep mean-field description closed while still producing nontrivial, channel-resolved dynamics; it enforces enough heterogeneity across neurons that the degenerate ``everyone receives the same effective update'' modes emphasized in Nguyen \& Pham (2023) are not the right mental model for the architecture we analyze.

Concretely, the non-verbatim mathematical content relative to Nguyen \& Pham (2023) is not a single lemma in isolation but the redefinition of the forward and backward objects: the channelized map
\[
  H_\ell=\sum_{k=1}^r W^{(\ell)}_{\cdot,k}\,f_k^{(\ell)}
\]
propagates discrepancies through $\|W^{(\ell)}\|_{\infty,1}$ and $\max_{k\le r}$ over channels (Bi-Lipschitz / coupling chain in our manuscript appendix). Frozen first-layer features also remove one delicate homotopy step used in full-rank analyses to keep first-layer support from collapsing during training, because dense-span random-feature richness can be enforced at initialization and preserved by freezing $W^0$.

We will revise the introduction and discussion to make the target domain explicit. We do not claim to subsume contemporary end-to-end training of billion-parameter models; we aim to contribute a theoretically grounded route to low-rank / structured training that is especially relevant when oscillatory or multiscale structure matters (scientific machine learning, PDE learning, and other settings where aggressively orthogonalized optimizers and full-rank dynamics are not the only relevant reference class). In that sense, the ``restriction'' is also an assumption that matches a family of applications where random features and factorized weights are already standard engineering tools.

For Theorem~4.1, we will keep the two-point, two-channel calculation as the fully proved anchor, and we will add a self-contained remark (feature-learning mechanism; see Q4 below) explaining the positive feedback through $H_2=\sum_k W_{c_2,k} f_k$ and ReLU gating at the ODE level. High-dimensional, overlapping features will be discussed honestly as primarily empirical and as ongoing work, not as something the current theorem already covers.

Q2 (Symmetry sensitivity; exponential factor in finite-width bound)

We take this concern seriously because it touches both what the theorems mean and what the experiments show.

On the Gr\"onwall exponential and ``symmetry loss.'' The exponential prefactor in Theorem~3.3-style bounds is the standard price of a worst-case stability analysis: in the quantitative coupling argument one obtains a constant of the schematic form $C_{\exp}=\exp\bigl(K_T(1+rK)\bigr)$ multiplying discretization and $1/\sqrt{n}$ errors. This is a conservative mathematical artifact. It should not be read as a causal explanation of a particular empirical symmetry diagnostic, and it is not something we use to claim that ``rank must be tiny or the mean-field picture breaks.'' In practice, symmetry-like structure can emerge (or disappear) from optimization noise, batching, initializations, and the target function, largely in parallel with whatever worst-case coupling constant appears in a bound.

On rank $10$ vs.\ $20$ and robustness. We agree that any phenomenon that exists only for a razor-thin set of hyperparameters is not a convincing ``mechanism.'' We also want to avoid dueling anecdotes: different targets, batch sizes, and initializations can change qualitative symmetry observations, and small-batch SGD already supplies many reasons for symmetry breaking even when the architecture admits symmetric solutions.

To make this concrete, we ran a focused reproducibility sweep after the review (fixed hidden width $1024$, varying rank, multiple seeds) and measured an even-function symmetry diagnostic on a synthetic even target (sum of cosines). The table below reports mean$\pm$std over five seeds ($42$--$46$). The quantitative story is not monotonic in the way one might guess from a single pilot run: in this protocol, $r=20$ produced lower mean symmetry error than $r=10$, while test MSE was comparable within noise. This does not refute your experience---it underscores that the symmetry observable is setup-dependent---but it motivates us to report variance and to tone down any universal claim about ``rank $10$ good, rank $20$ bad.''

Table (post-review symmetry diagnostic sweep; synthetic even target; MMNN with frozen features; width $1024$). $D_{\mathrm{sym}}=\mathbb{E}\bigl[(f(x)-f(-x))^2\bigr]$ on a positive $x$ grid; ``rel.'' divides by mean output energy. Exact replication details will be in supplementary material.

\begin{center}
\begin{tabular}{r|l|l|l}
Rank $r$ & test MSE (mean$\pm$std) & $D_{\mathrm{sym}}$ (mean$\pm$std) & rel.\ $D_{\mathrm{sym}}$ (mean$\pm$std) \\
\hline
$10$ & $1.10\pm0.26 \times 10^{-3}$ & $1.89\pm1.73 \times 10^{-4}$ & $1.57\pm1.42 \times 10^{-4}$ \\
$20$ & $1.00\pm0.22 \times 10^{-3}$ & $4.18\pm0.58 \times 10^{-5}$ & $3.49\pm0.48 \times 10^{-5}$
\end{tabular}
\end{center}

We will fold a shorter version of this discussion into the camera-ready (or supplement), and we will phrase symmetry observations as hypotheses supported by particular protocols rather than as necessary consequences of low rank.

Q3 (Scaling to higher intrinsic dimension; phase transition)

We agree this is one of the most important ``beyond the theorem'' questions. Our current mean-field guarantees are formulated for fixed rank $r$ in the infinite-width limit; they are not, by themselves, a phase diagram in $(d,r,N)$. We will state that limitation plainly in the discussion.

That said, we can still articulate a research roadmap that connects to work in progress. On the NTK / kernel side, understanding how rank interacts with input dimension requires controlling the induced kernel ensemble; related random-matrix analyses become delicate for low-rank structured weights, and we are actively building on recent frameworks for structured kernels (e.g.\ \url{https://arxiv.org/abs/2508.20036}) to connect $r$ scaling to high-dimensional concentration. On the mean-field / Chizat side, recent scaling-law and mean-field scaling discussions suggest that a ``sweet spot'' for faithful ODE descriptions can occur when rank grows sublinearly with width (informally, regimes such as $r\sim\sqrt{\mathrm{width}}$ are often discussed as practically relevant; see \url{https://arxiv.org/pdf/2509.10167} for related scaling perspectives). Separately, in LoRA-style NTK analyses, one sometimes sees rank thresholds of the form $r\gtrsim M^{\alpha}$ (with $\alpha\in[1/4,1/2]$) discussed as regimes where kernel descriptions remain predictive; we treat such statements as indicators, not as theorems proved for our exact architecture.

Regarding the specific ``exponential suppression'' intuition: in the ReLU gating argument, the problematic configuration corresponds to a joint sign-alignment event across channels. The heuristic is not that optimization fails globally once $r$ grows, but that the input-region measure where all channels are simultaneously dead can shrink exponentially in $r$ under symmetric mixing assumptions. In higher intrinsic dimension, whether this remains effective depends on how fast $r$ must grow to represent the target versus how concentration/gating geometry scales with $d$; this is exactly why we frame $(d,r,N)$ as an empirical+theoretical phase-diagram question rather than a settled theorem.

Empirically, we have begun post-review sweeps on synthetic high-dimensional targets (using mixed sampling to avoid degenerate targets in high $d$) and we will report these curves without over-claiming a sharp phase transition unless the data clearly supports one. Any clean $(d,r)$ transition is, at present, conjectural and belongs in ``future work'' unless we can prove it.

Q4 (Mechanism behind the channel spike --- beyond the two-point toy)

The explanatory content is not the simplified two-point, two-channel statement itself, but the positive feedback it isolates.

Fix a second-layer neuron index $c_2$ and recall the low-rank pre-activation
\[
  H_2(c_2;x,W)=\sum_{j=1}^r W_{c_2,j}\,f_j(x;W).
\]
For ReLU, gating enters mean-field gradients through $\varphi'_2(H_2)=\mathbf{1}\{H_2>0\}$. The channel-$k$ backward signal (cf.\ the ODE for $\partial_t w_1(\cdot,k)$) is built from expectations over $C_2$ of terms that couple (i) the mixing weight $W_{C_2,k}$, (ii) the gate $\mathbf{1}\{H_2(C_2;x,W)>0\}$, and (iii) upstream factors (output error and later-layer weights).

When channel $k$ dominates, those neurons $c_2$ that contribute to the channel-$k$ drift are precisely those with the gate ``on.'' Across $c_2$, the effective weights in the expectation are then positively aligned with $W_{c_2,k}$ and with $\mathbf{1}\{H_2(c_2)>0\}$, producing a large contribution to the drift of the $k$-th channel, which in turn pushes $f_k$ further in the same direction: a reinforcing loop. The ODE-level mechanism is dimension-agnostic; the two-point theorem is a minimal tractable instance.

Q5 (Limitations: frozen mixing as an architectural restriction)

We agree it would be misleading to suggest that every practical deep network is well modeled by completely frozen mixing. At the same time, we want to articulate the positive claim we are actually making: random features and factorized weights are not exotic---they are standard tools in scientific computing and structured learning---and the point of our analysis is that global optimality-type conclusions can still be compatible with substantial low-rank structure in a mean-field scaling, contrary to a casual intuition that ``low rank only makes optimization worse.'' We will rewrite the limitations paragraph so it reads as a scope statement (what the theorems assume and what they imply within that class) rather than as an apology that could be mistaken for admitting the model is irrelevant to modern ML. We will also draw a clearer contrast with optimization stacks that deliberately orthogonalize updates (e.g.\ some large-scale training recipes), which target a different design point than the PDE / SciML settings we emphasize.
